\documentclass[a4paper,12pt]{article}
\setlength{\headheight}{68.803pt}
\addtolength{\topmargin}{-56.803pt}

\usepackage[utf8]{inputenc}
\usepackage[T1]{fontenc}
\usepackage{lmodern}
\usepackage[nopatch=footnote]{microtype}
\usepackage{fvextra}
\DefineVerbatimEnvironment{Verbatim}{Verbatim}{
  breaklines=true,
  breakanywhere=true,
  frame=single,
  fontsize=\small
}

\usepackage[
  left=1in,
  right=1in,
  top=3cm,
  bottom=1in
]{geometry}

\usepackage[
    colorlinks=true,
    linkcolor=black,
    urlcolor=blue,
    citecolor=blue,
    pdfborder={0 0 0},
    pdfauthor={Massimo Mantineo},
    pdftitle={MidTerm HO14: Sistema Bancario TCP},
    pdfsubject={Laboratorio di Reti e Sistemi Distribuiti: HANDS-ON 14},
]{hyperref}

\usepackage{graphicx}
\graphicspath{{../../assets/}{../../assets/ho1/}{../../assets/ho2/}{../../assets/ho3/}{../../assets/ho4/}{../../assets/ho5/}{../../assets/ho6/}{../../assets/ho7/}{../../assets/ho8/}{../../assets/ho10/}}
\usepackage{tikz}
\usetikzlibrary{positioning, arrows.meta, shapes.symbols}
\usepackage{listings}
\usepackage{xcolor}
\usepackage{amsmath, amssymb}
\usepackage{fancyhdr}
\usepackage{setspace}
\usepackage{pgfplots}
\pgfplotsset{compat=1.17}
\usepackage{booktabs}
\usepackage[skins, listings]{tcolorbox}
\usepackage{float}

\newtcblisting{terminal}[1][]{
    listing only,
    colback=black!85!white,
    colframe=black!70!white,
    colupper=green!80!white,
    coltitle=white,
    title=#1,
    fonttitle=\bfseries\sffamily,
    listing options={
        language={},
        basicstyle=\ttfamily\small\color{green!80!white},
        backgroundcolor={},
        frame=none,
        breaklines=true,
        showstringspaces=false,
        numbers=none
    },
    enhanced,
    drop shadow,
    arc=3mm
}

\newtcblisting{ccode}[1][]{
    listing only,
    colback=blue!3!white,
    colframe=blue!70!black,
    title=#1,
    fonttitle=\bfseries\sffamily,
    listing options={
        style=cstyle,
        backgroundcolor={},
        frame=none
    },
    enhanced,
    drop shadow,
    arc=3mm
}

\setlength{\footskip}{50pt}

\lstdefinestyle{cstyle}{
    language=C,
    basicstyle=\ttfamily\small,
    numbers=left,
    numberstyle=\tiny,
    stepnumber=1,
    numbersep=5pt,
    backgroundcolor=\color{white},
    showspaces=false,
    showstringspaces=false,
    showtabs=false,
    frame=single,
    tabsize=4,
    captionpos=b,
    breaklines=true,
    breakatwhitespace=true,
    keywordstyle=\color{blue}\bfseries,
    commentstyle=\color{green!50!black},
    stringstyle=\color{red},
}
\lstset{style=cstyle}

\pagestyle{fancy}
\fancyhf{}

\fancyhead[L]{%
  \parbox[b]{0.45\textwidth}{%
    \small
    Student Name: Massimo Mantineo\\
    Student ID: 541924
    \vspace{20pt}
  }%
}
\fancyhead[C]{%
  \parbox[b]{0.2\textwidth}{%
    \centering
    \normalsize \bfseries
    LabReSiD25\\
    Hands-On 14
    \vspace{20pt}
  }%
}
\fancyhead[R]{%
  \parbox[b]{0.35\textwidth}{%
    \flushright
    \includegraphics[width=3.5cm]{\detokenize{logo_unime_orizontale.png}}
    \vspace{15pt}
  }%
}

\renewcommand{\contentsname}{Indice}

\title{\vspace{-2cm}Report MidTerm (HO14)\\
\large Sistema Bancario Multithread TCP}
\author{Massimo Mantineo \\ Matricola: 541924}
\date{MESSINA, A.A. 2024/2025}

\begin{document}

\begin{titlepage}
    \thispagestyle{empty}
    \centering
    \vspace*{1cm}
    \includegraphics[width=6cm]{\detokenize{logo_unime_orizontale.png}}\par\vspace{1cm}
    \vspace{2cm}
    {\Large \textbf{Università degli Studi di Messina}\par}
    {\large Dipartimento MIFT - Corso di Laurea Triennale in Informatica (L-31)\par}
    \vspace*{1cm}
    {\large Laboratorio di Reti e Sistemi Distribuiti\par}
    \vspace{1cm}
    {\Large \textbf{HANDS-ON 14 (MidTerm):}\\
    Sviluppo di un Sistema Bancario TCP Multithread\par}
    \vspace{2cm}
    {\large Massimo Mantineo \par}
    {\large Matricola: 541924 \par}
    \vfill
    {\large Messina, A.A. 2024/2025 \par}
\end{titlepage}

\clearpage
\pagestyle{fancy}
\tableofcontents
\newpage

\section{Problem Statement}
Il MidTerm (Hands-On 14) riunisce ed eleva le competenze acquisite operando finora con i Socket Berckley e con le astrazioni della libreria \texttt{Pthreads}. La richiesta principale è l'implementazione in linguaggio C di un \textbf{Sistema Bancario} minimale, in grado di gestire un correntista che invia comandi CRUD asincroni da molteplici istanti d'accesso (Terminali differenti, App, E-banking sportello).

\begin{figure}[H]
    \centering
    \begin{tikzpicture}[
        node distance=1.2cm,
        client/.style={rectangle, draw=blue!70, fill=blue!5, thick, minimum width=2.5cm, minimum height=0.8cm, font=\small\sffamily},
        server/.style={rectangle, draw=red!70, fill=red!5, thick, minimum width=4cm, minimum height=1.5cm, align=center, font=\bfseries\sffamily},
        db/.style={cylinder, draw=orange!70, fill=orange!5, thick, shape border rotate=90, minimum width=2cm, minimum height=2cm, font=\small\sffamily},
        arrow/.style={-Stealth, thick, gray!80}
    ]
        % Clients
        \node[client] (c1) {Terminale};
        \node[client, below=0.3cm of c1] (c2) {App Mobile};
        \node[client, below=0.3cm of c2] (c3) {Web Banking};

        % Server
        \node[server, right=3cm of c2] (serv) {BANCOMAT SERVER\\(TCP Multithreaded)};

        % Database
        \node[db, right=3cm of serv] (database) {Account\\Balances};

        % Connections
        \draw[arrow] (c1.east) -- (serv.north west);
        \draw[arrow] (c2.east) -- (serv.west);
        \draw[arrow] (c3.east) -- (serv.south west);
        
        \draw[arrow, <->, orange!70, very thick] (serv) -- node[above, font=\tiny, color=black] {Thread-Safe Access} (database);
        
        \node[below=0.1cm of serv, font=\itshape\tiny] {Mutex Protected Resources};
    \end{tikzpicture}
    \caption{Architettura del Sistema Bancario Distribuito Multi-Client.}
    \label{fig:banking_arch}
\end{figure}

Il sistema si fonda su un'architettura **Client-Server TCP** e prevede limitatori rigorosi per la concorrenza (10 Threads), più protezioni di memoria contro la \textit{Race Condition} che andrebbe a distruggere la congruenza dell'estratto conto.

\section{Stato dell'Arte}
Architetture come quella proposta si pongono alla base dei moderni layer di accesso Cloud.
\begin{itemize}
    \item \textbf{TCP Connection-Oriented}: A differenza dell'UDP (Connectionless), la TCP Handshake assicura che il traffico bancario sia controllato, privo di perdita pacchetti (\textit{Acknowledgment}) ed in corretto ordine seriale.
    \item \textbf{Gestione Demultiplexer-Worker}: Il design pattern adottato è quello \textit{One-Thread-Per-Connection}. Il \texttt{main()} accetta pacchetti SYN, istanzia la socket e ne \textit{esporta} il File Descriptor allocando un Worker Thread figlio staccato (Detached).
    \item \textbf{Critical Section}: Le operazioni di bilanciamento conto (Lettura-Modifica-Scrittura) non sono atomiche ai cicli di clock. Il DB (simulato) sarà interamente protetto tramite il costrutto \texttt{pthread\_mutex\_t}.
\end{itemize}

\section{Metodologia ed Implementazione}

\subsection{Progettazione del Database}
La struttura \texttt{Transaction} immagazzina l'ID, la causale e l'importo. Un array \texttt{db[MAX\_RECORDS]} simula la tabella centrale di archiviazione in backend. La discriminante principale è lo stato logico di esistenza.
\begin{ccode}[Struttura Transazione]
typedef struct {
    int id;
    float amount;
    char reason[256];
    int active;
} Transaction;
\end{ccode}

\subsection{Protezione Logica e Limitazione Server}
Per obbedire al vincolo "Al massimo 10 thread", abbiamo implementato un semaforo limitante pre-accept:
\begin{ccode}[Limitatore di traffico Server]
// Inizializzato a MAX\_NUM (10)
sem\_wait(&thread\_limiter); 
client\_sock = accept(server\_sock, ...);
// Viene passata la richiesta a pthread\_create
\end{ccode}
Ogniqualvolta una richiesta accede al DB, acquisisce un \texttt{pthread\_mutex\_lock(\&db\_lock)}. Se è un inserzione (\texttt{ADD}), scansiona i duplicati di ID per prevenire collisioni logiche. Lo stream testuale inviato dai socket viene parsate comodamente tramite \texttt{sscanf}.


\subsection{Istruzioni di Esecuzione}
L'applicativo viene gestito interamente tramite Makefile:
\begin{terminal}[Build & Run]
$ make
$ ./programma
\end{terminal}

\section{Risultati}
Innescando il demone server ed aprendo il client tramite l'eseguibile \texttt{./client}, la suite risponde correttamente con l'estratto CRUD implementato:
\begin{terminal}[Esecuzione Client e Server RPC]
massimo@PORTATILE-Linux:~/Scrivania/LabReSid24-25/hands-on/ho14$ ./client
--- Connected to Bank Server ---
Available commands:
  ADD <id> <amount> <reason>
  UPDATE <id> <amount> <reason>
  DELETE <id>
  LIST
  QUIT
BANK> ADD 1 1500.00 Stipendio Mese Marzo
OK: ADDED 1
BANK> ADD 2 55.40 Rimborso Pizza
OK: ADDED 2
BANK> LIST
--- TRANSACTION LIST ---
[1] 1500.00 - Stipendio Mese Marzo
[2] 55.40 - Rimborso Pizza
------------------------
BANK> UPDATE 2 -20.00 Storno Parziale
OK: UPDATED 2
BANK> DELETE 1
OK: DELETED 1
BANK> LIST
--- TRANSACTION LIST ---
[2] -20.00 - Storno Parziale
------------------------
BANK> QUIT
\end{terminal}
L'architettura garantisce una stabilità duratura. La chiusura dell'innesco `client` deallocherà silenziosamente il file descriptor riconsegnando un token libero al \texttt{thread\_limiter}.

\subsection{Analisi di Rete (Wireshark)}
I test condotti sull'interfaccia di Loopback confermano a livello ISO/OSI il corretto funzionamento dell'inizializzazione Client/Server.

\begin{figure}[H]
    \centering
    \includegraphics[width=1\textwidth]{\detokenize{ho14/ho14_wireshark_rpc.png}}
    \caption{Cattura Wireshark: TCP 3-Way Handshake, Invio Dati (PSH) e Teardown}
\end{figure}

Come richiesto dal tracciato teorico, si osservano chiaramente le tre iterazioni fondamentali gestite dallo strato TCP:
1. L'avvio con l'Handshake a tre vie (\texttt{[SYN]}, \texttt{[SYN, ACK]} da parte del server, e l'ultimo \texttt{[ACK]}) che instaura il tunnel virtuale sicuro.
2. L'invio bidirezionale del payload mediante pacchetti identificati dal flag \texttt{[PSH, ACK]}, dove si cela l'operazione \texttt{ADD} inviata in chiaro.
3. Lo smantellamento concordato alla pressione del \texttt{QUIT} pilotato dal flag \texttt{[FIN, ACK]}.


\section{Conclusioni e Sviluppi Futuri}
A un livello di Proof-Of-Concept, la gestione è completamente coerente: non permette transazioni duplici, non permette cancellazioni inesatte e parallelizza correttamente client multipli.

\noindent \textbf{Sviluppi futuri ideali:}
\begin{itemize}
    \item \textbf{DB Storage Persistente}: Abbandono dell'array per delegare salvataggi espliciti su file \texttt{.txt} o database relazionali veri e propri (PostgreSQL, SQLite), così il riavvio del server non annienta le operazioni precedenti.
    \item \textbf{Read-Write Locks}: Utilizzo di \texttt{pthread\_rwlock\_t} al posto del grezzo mutex per far sì che molteplici chiamate \texttt{LIST} possano leggere senza bloccarsi vicendevolmente nel caso di estremo traffico in sola visualizzazione.
\end{itemize}

\end{document}
